\section{\raggedright Particle and generation assignments}
The algebraic carrier suggests several representation-theoretic
ingredients: unit quaternions provide an $SU(2)$ action, complex scalars
provide phase transformations, and the octonionic setting motivates
colour representations. These ingredients are useful starting points,
but they do not uniquely fix the physical gauge theory. In particular,
a common complex phase is not by itself the electromagnetic charge
operator or a derivation of hypercharge.

The distinction between colour and generation is equally important.
A three-dimensional subspace carrying a colour action cannot be
identified with three generations merely by counting its basis vectors.
The Companion's inherited slot assignments require additional choices
of embeddings, charge operators and the relation between internal
representations and physical states. We therefore retain the gauge and
generation discussion as a conditional assignment problem, not as a
derivation of the complete Standard Model.

A more explicit comparison construction starts from a seed isometry
$T_0:\mathbb C^3\to\mathcal K$ and an orthogonal closure projector
$P_c$ on a declared common carrier. If $T_0^*P_cT_0$ is positive
definite, then
\begin{equation}
 S=P_cT_0(T_0^*P_cT_0)^{-1/2},\qquad S^*S=I_3.
\end{equation}
This constructs three orthonormal closed states under its hypotheses.
It does not determine the seed embedding or explain why those states
are lepton generations. We distinguish $P_c$ from the HC projector $C$:
closure within a selected space and even/odd HC grading are separate
structures. This distinction becomes necessary when a coupling is
used to define mass readout.
\subsection{What a representation construction must specify}
The algebraic question is not whether symbols resembling the three
gauge factors can be placed in a diagram. It is which operators act
on which states, how the actions commute, and which observables are
assigned to their invariant subspaces. A physical gauge interpretation
also requires local transformation laws, connections and the appropriate
matter representations. The compact groups encountered in a carrier
are only part of these data.

For the quaternion factor, write the Hermitian generators as
$T_k=\sigma_k/2$. The Pauli multiplication law gives
\begin{equation}
 [T_j,T_k]=i\epsilon_{jkl}T_l.
\end{equation}
This exhibits the Lie algebra of the unit-quaternion $SU(2)$ action.
It does not decide whether a particular subspace represents a weak
doublet, a spatial spinor or another two-component system. That
decision depends on the remaining representation and coupling data.
Likewise the commuting scalar phase $e^{i\theta}$ acts uniformly on
the carrier until a nontrivial charge operator is specified.

The octonionic ingredient is equally sensitive to the chosen complex
structure. Fix an imaginary unit $e_7$ and let $I_7=L_{e_7}$ denote
left multiplication on the real octonion space. Here $I_7^2=-I$,
so it is a complex structure, not the involution $J$ used for the HC
grading. Relative to it, the real eight-dimensional space can be
read as a complex four-dimensional space. The familiar singlet/triplet
organization associated with the stabilizer of $e_7$ motivates the
colour comparison used in the Companion. All identifications in this
paragraph refer to that selected internal complex structure.

On the complexification, the corresponding complex-linear projections
are
\begin{equation}
 \Pi_\pm=\frac{I\mp iI_7}{2},\qquad
 \Pi_\pm I_7=\pm i\Pi_\pm.
\end{equation}
Complex coefficient conjugation exchanges these projections. They
must not be substituted for $C$ or for a closure projector without
a map that identifies the relevant domains and meanings. The same
letter ``projector'' can otherwise hide three different decompositions:
internal complex structure, HC parity and physical closure.

\subsection{A conditional ladder construction and its exact counting}
The charge construction retained in the Companion can be separated
into an algebraic hypothesis and the consequences that follow from it.
In its specified octonion multiplication convention, introduce
\begin{equation}
 \alpha_a=\tfrac12(L_{e_a}+iL_{e_7e_a}),\quad a=4,5,6,
 \qquad N=\sum_{a=4}^6\alpha_a^*\alpha_a,
 \qquad Q_{\rm ch}=N/3.
\end{equation}
The represented Clifford multiplication and adjoint must give the
canonical anticommutation relations
\begin{equation}
 \alpha_a\alpha_b+\alpha_b\alpha_a=0,\qquad
 \alpha_a\alpha_b^*+\alpha_b^*\alpha_a=\delta_{ab}I.
\end{equation}
Once these relations and a cyclic vacuum annihilated by the three
$\alpha_a$ are supplied, the occupation counting is elementary.
Each creation operator can occur at most once, because its square
vanishes. The states with $k$ occupied modes therefore have dimension
$\binom3k$, and $N$ has eigenvalue $k$ on that subspace. The charge
comparison consequently has the spectrum
\begin{equation}
 \begin{array}{c|rrrr}
 k&0&1&2&3\\\hline
 \text{multiplicity}&1&3&3&1\\
 Q_{\rm ch}&0&1/3&2/3&1
 \end{array}
\end{equation}
The total dimension is eight, the dimension of the three-mode fermionic
Fock representation. These multiplicities are consequences of the
chosen ladder algebra. They do not establish three lepton families.

An internal $SU(3)$ action that rotates the creation operators and
preserves the vacuum acts on the one-occupation space as a triplet,
and on the two-occupation space as its conjugate, since these are
the first and second exterior powers of a three-dimensional complex
space. Vacuum and full occupation are singlets. This gives a precise
representation question: is the represented colour action on the
proposed octonionic carrier the action required by these operators?
It must be checked with the chosen multiplication and adjoint; the
counting table alone is not that check.

The Companion's historical calculations motivate this compatibility,
but the current paper does not turn a preserved numerical calculation
into a new formal theorem. Its physical use still needs the embedding
of particles, their chirality and the electroweak charge assignments.
In particular a number operator with the displayed eigenvalues does
not, by itself, derive hypercharge or the full weak representation.

This decomposition also explains a correction to the literal slot
dictionary. With the inner complex structure, vectors such as
$e_4,e_5,e_6$ can form a complex triplet basis only because their
$I_7$-images supply the corresponding real partners. Counting the
images as another independent physical triplet double-counts the
same complex degrees of freedom. Changing to the commuting outer
complex unit changes the decomposition into charge eigenvectors.
State labels must follow that change; an unqualified equality of
charges across the proposed frames is not justified.

\subsection{The independent generation space}
Three colours within one representation and three generations of a
given particle are different multiplicities. If the colour action is
nontrivial on a three-dimensional subspace, that action does not
automatically act as a generation permutation. A generation assignment
needs a second declared structure and its relation to the carrier.
The closure-seed construction supplies a conditional way to express
that requirement without identifying raw octonion coordinates with
mass eigenstates.

Let $T_0:\mathbb C^3\to\mathcal K$ be an isometric seed embedding,
and let $P_c$ be an orthogonal projector. Define
$G_0=T_0^*P_cT_0$. If $G_0>0$, then the states in
\begin{equation}
 S=P_cT_0G_0^{-1/2}
\end{equation}
are closed and orthonormal:
\begin{equation}
 P_cS=S,\qquad
 S^*S=G_0^{-1/2}T_0^*P_c^2T_0G_0^{-1/2}=I_3.
\end{equation}
If $G_0$ has rank below three, projection has removed an independent
seed direction and cannot produce three independent states by this
formula. If $G_0$ is not diagonal, normalization already mixes the
original seed slots. Thus even a successful construction does not
identify each initial slot separately as a physical mass eigenstate.

Suppose also that $[P_c,J]=0$ and $JT_0=-T_0$. The same calculation
gives $JS=-S$. This realizes three closed odd inputs. A nonzero
operator restricted to this odd space preserves oddness and commutes
with the restricted $J=-I$. It cannot also anticommute with $J$.
An HC-odd coupling must connect opposite parity sectors; an operator
that merely has odd outputs is a different object.

The distinction supplies a useful diagnostic before computing any
spectrum. A proposed Dirac or mass coupling should specify whether
it acts within a sector, changes HC grading, changes closure sector,
or combines these operations. Its matrix blocks and domains determine
the answer. Calling the operator ``odd'' or giving it a familiar
letter does not determine which of these statements is intended.

\subsection{A represented transition and the choice of mass basis}
Take a bounded dimensionless comparison pacer $\widetilde R$ on the
same carrier and define
\begin{equation}
 D=(I-P_c)\widetilde R P_c,\qquad B=DS,\qquad Q=B^*B.
\end{equation}
Then $Q\geq0$. If a mass law $M_{\rm mass}=\kappa Q$ is independently
adopted, its candidate eigenmasses are $\kappa q_i$, where $q_i$
are the eigenvalues of $Q$. A fixed-slot readout instead uses
$\kappa Q_{ii}=\kappa\|B\varepsilon_i\|^2$. These prescriptions
need not coincide.

For example, with $\boldsymbol1=(1,1,1)^{\mathsf T}$, let
\begin{equation}
 Q=\tfrac13\boldsymbol1\boldsymbol1^*.
\end{equation}
Its three diagonal entries equal $1/3$, but its spectrum is $(1,0,0)$.
The fixed-slot mass-root ratio would be $1/3$, whereas the spectral
ratio would be $1$. These are different observables of one positive
matrix. The example excludes the inference from equal diagonal
couplings to degenerate physical eigenmasses.

There is also a rank requirement. Since
$\langle v,Qv\rangle=\|Bv\|^2$, the kernels of $Q$ and $B$
coincide, and
\begin{equation}
 \operatorname{rank}Q=\operatorname{rank}B.
\end{equation}
Three strictly positive candidate masses require three independent
transition outputs. A single scalar output channel has rank at most
one, regardless of how many input labels it is given. An operator-valued
channel may have higher dimension, but that dimension must be represented
rather than inferred from the name of a probe channel.

\subsection{What mixing angles would require}
Mixing is a relative-basis observable. To make the requirement explicit,
consider two Hermitian generation comparisons $Q_1,Q_2$ with specified
physical meanings. Let $V_1,V_2$ diagonalize them in the same generation
space. Their relative unitary matrix is
\begin{equation}
 U_{12}=V_1^*V_2.
\end{equation}
Changing the common generation coordinates by $F$ replaces the two
diagonalizers by $FV_1,FV_2$ and leaves $U_{12}$ unchanged. Rephasing,
ordering and eigenvalue degeneracies require their usual separate
conventions before individual parameters can be interpreted.

This construction identifies the data a mixing prediction would need:
two independently represented operators, their common kinetic metric,
their assignment to physical species, and a rule fixing their relative
eigenspaces. A common cyclic action or the dimension three does not
provide those data. In a degenerate eigenspace the basis is free, so
apparent mixing angles obtained from an arbitrary diagonalizer are
not predictions. The mass-readout analysis below therefore treats
the spectrum, its physical scale and its generation basis as distinct
questions, rather than allowing one successful matrix pattern to
stand in for all three.
